% ============================================================
%  FORMAL SECTIONS — to integrate into the manuscript
%  Target journal: Modelling (MDPI), ISSN 2673-3951
%  Paper: "Optimizing pallet configuration for air transport
%          of cut flowers"
% ============================================================

% ============================================================
%  REQUIRED PACKAGES (add to preamble)
% ============================================================
% \usepackage{booktabs}
% \usepackage{amsmath, amssymb}
% \usepackage[ruled, vlined, linesnumbered]{algorithm2e}
% \usepackage{multirow}
% \usepackage{array}


% ============================================================
%  TABLE OF NOTATION
%  Place at the beginning of Section 4 (Problem Definition)
% ============================================================

\begin{table}[ht]
\caption{Mathematical notation used throughout this paper.}
\label{tab:notation}
\centering
\small
\begin{tabular}{@{}ll@{}}
\toprule
\textbf{Symbol} & \textbf{Description} \\
\midrule
\multicolumn{2}{@{}l}{\textit{Sets and indices}} \\
$N$               & Total number of boxes assigned to a single flight \\
$M$               & Total number of ULDs available on the aircraft \\
$\mathcal{B} = (b_1, b_2, \ldots, b_N)$ & Ordered sequence of boxes in arrival order \\
$\mathcal{U} = \{U_1, U_2, \ldots, U_M\}$ & Set of ULDs, indexed by their loading order \\
$\mathcal{R}(b_i)$ & Set of admissible axis-aligned rotations of box $b_i$ \\
\midrule
\multicolumn{2}{@{}l}{\textit{Box parameters}} \\
$(w_i, h_i, d_i)$ & Width, height, and depth of box $b_i$ (original orientation) \\
$(w_i', h_i', d_i')$ & Effective dimensions of $b_i$ after applying rotation $r_i \in \mathcal{R}(b_i)$ \\
$m_i$             & Gross weight of box $b_i$ (kg) \\
$\mathrm{vol}(b_i)$ & Volume of box $b_i$: $\mathrm{vol}(b_i) = w_i \cdot h_i \cdot d_i$ \\
\midrule
\multicolumn{2}{@{}l}{\textit{ULD parameters}} \\
$(W_j, H_j, D_j)$ & Internal width, height, and depth of ULD $U_j$ \\
$C_j^w$           & Maximum weight capacity of ULD $U_j$ (kg) \\
$\mathrm{Vol}(U_j)$ & Internal volume of ULD $U_j$: $\mathrm{Vol}(U_j) = W_j \cdot H_j \cdot D_j$ \\
\midrule
\multicolumn{2}{@{}l}{\textit{Decision variables}} \\
$x_{ij} \in \{0,1\}$ & 1 if box $b_i$ is assigned to ULD $U_j$; 0 otherwise \\
$r_i \in \mathcal{R}(b_i)$ & Selected rotation for box $b_i$ \\
$(p_i^x, p_i^y, p_i^z) \in \mathbb{N}^3$ & Placement origin (lower-left-back corner) of $b_i$ within its ULD \\
\midrule
\multicolumn{2}{@{}l}{\textit{State representation}} \\
$\mathcal{F}_j$   & Set of free (unoccupied) cuboid subspaces within ULD $U_j$ \\
$\mathbf{H}_j \in \mathbb{N}^{D_j \times W_j}$ & Height map of ULD $U_j$: $\mathbf{H}_j(z, x)$ gives the stacking height at column $(x,z)$ \\
$s = (s^x, s^y, s^z, s^W, s^H, s^D)$ & A free cuboid (split) with origin and dimensions \\
\midrule
\multicolumn{2}{@{}l}{\textit{Algorithm parameters}} \\
$k$               & Lookahead window size (number of observable upcoming boxes) \\
$\alpha$          & Minimum support fraction for stability constraint ($\alpha = 0.5$) \\
$\mathcal{H}$     & Heuristic policy: $\mathcal{H} \in \{\text{BL}, \text{BAF}, \text{BLSF}\}$ \\
\bottomrule
\end{tabular}
\end{table}


% ============================================================
%  SECTION 4.2 — MATHEMATICAL PROBLEM FORMULATION
%  Insert after existing Section 4.1 (Problem definition)
% ============================================================

\subsection{Mathematical Problem Formulation}
\label{sec:formulation}

The palletization problem studied in this work is formulated as an
\emph{online three-dimensional bin packing problem} (O-3D-BPP) over a
heterogeneous set of ULDs.
Table~\ref{tab:notation} summarises the notation used throughout.

\subsubsection*{Objective function}

The primary objective is to maximise the mean volumetric utilisation
across all $M$ ULDs on the aircraft:

\begin{equation}
\label{eq:objective}
\maximise \quad
\bar{U}_{\mathrm{vol}}
= \frac{1}{M} \sum_{j=1}^{M}
  \frac{\displaystyle\sum_{i=1}^{N} x_{ij}\, w_i' h_i' d_i'}
       {W_j H_j D_j},
\end{equation}

\noindent
where $w_i' h_i' d_i'$ denotes the volume of box $b_i$ under its
selected rotation $r_i$.
Equivalently, $\bar{U}_{\mathrm{vol}}$ measures the average fraction of
each ULD's internal volume that is occupied by cargo.

\subsubsection*{Constraints}

The optimisation is subject to the following constraints.

\paragraph{C1 — Unique assignment.}
Each box is assigned to at most one ULD:
\begin{equation}
\label{eq:unique}
\sum_{j=1}^{M} x_{ij} \leq 1, \quad \forall\, i \in \{1,\ldots,N\}.
\end{equation}

\paragraph{C2 — Geometric containment.}
If box $b_i$ is placed in ULD $U_j$, its entire volume must lie within
the ULD's interior:
\begin{equation}
\label{eq:containment}
x_{ij} = 1 \;\Rightarrow\;
\begin{cases}
  p_i^x + w_i' \leq W_j,\\
  p_i^y + h_i' \leq H_j,\\
  p_i^z + d_i' \leq D_j.
\end{cases}
\end{equation}

\paragraph{C3 — Non-overlap.}
No two boxes assigned to the same ULD may occupy overlapping volumes.
For all pairs $(b_i, b_{i'})$ with $x_{ij} = x_{i'j} = 1$ and $i \neq i'$,
the placed cuboids must be disjoint:
\begin{equation}
\label{eq:nooverlap}
\exists\, \delta \in \{x, y, z\} \text{ such that }
p_i^{\delta} + d_i'^{(\delta)} \leq p_{i'}^{\delta}
\;\text{ or }\;
p_{i'}^{\delta} + d_{i'}'^{(\delta)} \leq p_i^{\delta},
\end{equation}
where $d_i'^{(x)} = w_i'$, $d_i'^{(y)} = h_i'$, $d_i'^{(z)} = d_i'$.

\paragraph{C4 — Gravitational stability.}
A placement at height $p_i^y$ is feasible only if the fraction of the
box's base area that is supported at that height meets a minimum threshold:
\begin{equation}
\label{eq:stability}
\frac{\bigl|\{(x,z) : p_i^x \leq x < p_i^x + w_i',\;
                      p_i^z \leq z < p_i^z + d_i',\;
                      \mathbf{H}_j(z,x) = p_i^y\}\bigr|}
     {w_i' \cdot d_i'}
\;\geq\; \alpha,
\end{equation}
with $\alpha = 0.50$ in this study.
The placement height is determined by the height map as
$p_i^y = \max_{(x,z) \in \text{footprint}}\mathbf{H}_j(z,x)$,
ensuring boxes rest on previously placed cargo or on the pallet floor.

\paragraph{C5 — ULD weight capacity.}
The total gross weight loaded into each ULD must not exceed its structural
limit:
\begin{equation}
\label{eq:weight}
\sum_{i=1}^{N} x_{ij}\, m_i \leq C_j^w, \quad \forall\, j \in \{1,\ldots,M\}.
\end{equation}

\paragraph{C6 — Online (irreversible) decisions.}
At decision step $t$, the packing agent has access only to the next $k$
boxes in the arrival sequence:
$\mathcal{W}_t = (b_t, b_{t+1}, \ldots, b_{\min(t+k-1, N)})$.
Once box $b_t$ is placed, its position cannot be revised.
This constraint formalises the real-world conveyor-based operation in which
boxes arrive continuously and placement decisions must be made on-the-fly.

\medskip
\noindent
Constraints C1--C5 are deterministic feasibility requirements enforced at
every step.
Constraint C6 restricts the information available to the decision maker,
rendering the problem intractable for exact offline solvers and motivating
the use of efficient, online heuristic policies.


% ============================================================
%  SECTION 5 — SOLUTION APPROACH (rewritten subsections)
%  Replace / expand the existing subsections
% ============================================================

% ----  5.1  Space Representation  ---------------------------

\subsection{Space Representation via Maximal Cuboid Partitioning}
\label{sec:space-rep}

The free space inside each ULD $U_j$ is represented as a dynamic set
$\mathcal{F}_j$ of non-overlapping, axis-aligned \emph{maximal free
cuboids} (splits).
At initialisation, $\mathcal{F}_j = \{\text{Cuboid}(0,0,0,W_j,H_j,D_j)\}$.
When a box is placed, every free cuboid that intersects the occupied region
is replaced by up to six maximal sub-cuboids, and any sub-cuboid that is
entirely contained within another is discarded (dominance pruning).
This representation, described formally in Algorithm~\ref{alg:partition},
avoids explicit voxelisation and remains computationally tractable for
the box dimensions encountered in cut-flower operations.

A two-dimensional \emph{height map} $\mathbf{H}_j \in \mathbb{N}^{D_j \times W_j}$
is maintained in parallel.
Entry $\mathbf{H}_j(z, x)$ stores the current stacking height at column
$(x, z)$ of the pallet floor.
After placing a box with footprint
$[p^x, p^x + w') \times [p^z, p^z + d')$ at height $p^y + h'$, the height
map is updated as:
\begin{equation}
\label{eq:hmap-update}
\mathbf{H}_j(z, x) \leftarrow
\max\!\bigl(\mathbf{H}_j(z, x),\; p^y + h'\bigr),
\quad (x,z) \in \text{footprint}.
\end{equation}
The height map is used both for candidate height computation
(Equation~\ref{eq:stability}) and as a compact state representation
consumed by the heuristic scoring functions.

% ----  5.2  Heuristic Placement Policies  -------------------

\subsection{Heuristic Placement Policies}
\label{sec:heuristics}

Three deterministic placement heuristics are evaluated.
Each heuristic defines a \emph{scoring function} over the set of
feasible placements $\mathcal{A}_t$ generated at step $t$
(see Section~\ref{sec:procedure}).
The placement with the minimum score is selected and executed
(ties are broken lexicographically on the remaining score components).

\subsubsection{Bottom-Left (BL)}
\label{sec:bl}

The BL heuristic selects the placement that minimises the resulting
occupied height of the loaded item, breaking ties first by the
rightward extent and then by the forward extent:
\begin{equation}
\label{eq:score-bl}
\sigma_{\mathrm{BL}}\bigl(b_i, r_i, (p^x, p^y, p^z)\bigr)
= \bigl(p^y + h_i',\; p^x + w_i',\; p^z + d_i'\bigr).
\end{equation}
The placement chosen is $\argmin_{\mathcal{A}_t} \sigma_{\mathrm{BL}}$
under lexicographic ordering.
By systematically minimising vertical growth, BL promotes compact,
layered packing patterns that are well suited to the rectangular
geometry of aircraft pallets.

\subsubsection{Best Area Fit (BAF)}
\label{sec:baf}

BAF selects the placement that \emph{minimises the volume of the hosting
free cuboid} $s$, preferring the tightest available space.
Among placements that share the same host volume, the secondary criterion
minimises the residual along the shorter horizontal dimension:
\begin{equation}
\label{eq:score-baf}
\sigma_{\mathrm{BAF}}\bigl(b_i, r_i, s\bigr)
= \bigl(\mathrm{Vol}(s),\;
         \min(s^W - w_i',\; s^H - h_i')\bigr),
\end{equation}
where $s^W$ and $s^H$ are the width and height of the hosting free
cuboid.
BAF is a three-dimensional generalisation of the Best Area Fit strategy
\citep{Hassan2019}: by occupying the most constrained available space
first, it tends to preserve larger contiguous free regions for subsequent
items.

\subsubsection{Best Long Side Fit (BLSF)}
\label{sec:blsf}

BLSF selects the placement that minimises the residual along the
\emph{longer} horizontal dimension of the hosting free cuboid:
\begin{equation}
\label{eq:score-blsf}
\sigma_{\mathrm{BLSF}}\bigl(b_i, r_i, s\bigr)
= \bigl(\max(s^W - w_i',\; s^H - h_i')\bigr).
\end{equation}
This strategy aims to reduce fragmentation by aligning items against the
dominant boundary of each free space, preserving elongated free regions
that are easier to fill with subsequent items
\citep{Martinez2021}.

% ----  5.3  Overall Packing Procedure  ----------------------

\subsection{Multi-ULD Online Packing Procedure}
\label{sec:procedure}

Algorithm~\ref{alg:main} describes the complete sequential packing
procedure executed for a single charter flight.

\begin{algorithm}[ht]
\SetAlgoLined
\DontPrintSemicolon
\KwIn{%
  Box sequence $\mathcal{B} = (b_1, \ldots, b_N)$;
  ULD set $\mathcal{U} = (U_1, \ldots, U_M)$;
  heuristic $\mathcal{H}$;
  lookahead size $k$;
  support threshold $\alpha$}
\KwOut{Packing plan $\Pi$; volumetric utilisation $\bar{U}_{\mathrm{vol}}$}
\BlankLine
\tcp{Initialisation}
\ForEach{$U_j \in \mathcal{U}$}{
  $\mathcal{F}_j \leftarrow \{\mathrm{Cuboid}(0,0,0,W_j,H_j,D_j)\}$\;
  $\mathbf{H}_j \leftarrow \mathbf{0}^{D_j \times W_j}$\;
}
$j \leftarrow 1$; \quad $t \leftarrow 1$\;
\BlankLine
\tcp{Online packing loop}
\While{$t \leq N$ \textbf{and} $j \leq M$}{
  $\mathcal{W}_t \leftarrow (b_t, \ldots, b_{\min(t+k-1,N)})$
    \tcp*{observe lookahead window}
  $\mathcal{A}_t \leftarrow \textsc{FeasiblePlacements}(
      \mathcal{W}_t, U_j, \mathcal{F}_j, \mathbf{H}_j, \alpha)$\;
  \uIf{$\mathcal{A}_t = \emptyset$}{
    $j \leftarrow j + 1$
      \tcp*{current ULD is full; advance to next}
  }\Else{
    $(b_i^*, r_i^*, \mathbf{p}^*) \leftarrow \argmin_{\mathcal{A}_t}\,
        \sigma_{\mathcal{H}}$
      \tcp*{select via heuristic score}
    $c^* \leftarrow \mathrm{Cuboid}(p^{x*}, p^{y*}, p^{z*}, w_i'^*, h_i'^*, d_i'^*)$\;
    $\mathcal{F}_j \leftarrow \textsc{SpacePartition}(\mathcal{F}_j, c^*)$
      \tcp*{update free space}
    Update $\mathbf{H}_j$ via Eq.~(\ref{eq:hmap-update})\;
    Record $(b_i^*, U_j, r_i^*, \mathbf{p}^*)$ in $\Pi$\;
    Remove $b_i^*$ from sequence; \textbf{if} $i^* = t$ \textbf{then} $t \leftarrow t + 1$\;
  }
}
\BlankLine
$\bar{U}_{\mathrm{vol}} \leftarrow \frac{1}{M}\sum_{j=1}^{M}
  \frac{\sum_{(b_i,U_j,\cdot,\cdot)\in\Pi} \mathrm{vol}(b_i)}
       {\mathrm{Vol}(U_j)}$\;
\Return{$\Pi,\;\bar{U}_{\mathrm{vol}}$}\;
\caption{Multi-ULD Online 3D Bin Packing}
\label{alg:main}
\end{algorithm}

% ----  5.4  Feasible Placement Generation  ------------------

\medskip
Algorithm~\ref{alg:feasible} details how the set of feasible placements
$\mathcal{A}_t$ is constructed at each step.

\begin{algorithm}[ht]
\SetAlgoLined
\DontPrintSemicolon
\KwIn{%
  Lookahead items $\mathcal{W}_t$;
  active ULD $U_j$;
  free splits $\mathcal{F}_j$;
  height map $\mathbf{H}_j$;
  support threshold $\alpha$}
\KwOut{Set of feasible placements $\mathcal{A}_t$}
\BlankLine
$\mathcal{A}_t \leftarrow \emptyset$\;
\ForEach{$b_i \in \mathcal{W}_t$}{
  \ForEach{rotation $r \in \mathcal{R}(b_i)$}{
    $(w', h', d') \leftarrow \mathrm{dims}(b_i, r)$\;
    \ForEach{free split $s \in \mathcal{F}_j$ such that $s^W \geq w'$, $s^H \geq h'$, $s^D \geq d'$}{
      $(p^x, p^z) \leftarrow (s^x, s^z)$\;
      $p^y \leftarrow \max_{(x,z)\in[p^x, p^x+w')\times[p^z,p^z+d')}\mathbf{H}_j(z,x)$\;
      \tcp{Stability check (Eq.~\ref{eq:stability})}
      $\text{support} \leftarrow \tfrac{1}{w' d'}
          \bigl|\{(x,z):\mathbf{H}_j(z,x)=p^y,\;(x,z)\in\text{footprint}\}\bigr|$\;
      \If{$\text{support} \geq \alpha$
          \textbf{ and } $p^y + h' \leq H_j$}{
        $\mathcal{A}_t \leftarrow \mathcal{A}_t \cup \{(b_i, r, (p^x,p^y,p^z), s)\}$\;
      }
    }
  }
}
\Return{$\mathcal{A}_t$}\;
\caption{\textsc{FeasiblePlacements}: candidate placement generation}
\label{alg:feasible}
\end{algorithm}

% ----  5.5  Space Partitioning  -----------------------------

\medskip
Algorithm~\ref{alg:partition} describes the maximal-cuboid update used
in Algorithm~\ref{alg:main} after each placement.

\begin{algorithm}[ht]
\SetAlgoLined
\DontPrintSemicolon
\KwIn{Current free splits $\mathcal{F}$; placed cuboid $c$}
\KwOut{Updated free splits $\mathcal{F}'$}
\BlankLine
$\mathcal{F}' \leftarrow \emptyset$;\quad $\mathcal{F}_{\mathrm{new}} \leftarrow \emptyset$\;
\ForEach{$s \in \mathcal{F}$}{
  \uIf{$s$ intersects $c$}{
    $\mathcal{F}_{\mathrm{new}} \leftarrow \mathcal{F}_{\mathrm{new}}
        \cup \textsc{MaximalSplit}(s, c)$
      \tcp*{up to 6 sub-cuboids}
  }\Else{
    $\mathcal{F}' \leftarrow \mathcal{F}' \cup \{s\}$\;
  }
}
\tcp{Dominance pruning: discard sub-cuboids contained in another}
\ForEach{$s \in \mathcal{F}_{\mathrm{new}}$}{
  \If{$\nexists\, s' \in (\mathcal{F}' \cup \mathcal{F}_{\mathrm{new}})$
      such that $s' \neq s$ and $s'$ contains $s$}{
    $\mathcal{F}' \leftarrow \mathcal{F}' \cup \{s\}$\;
  }
}
\Return{$\mathcal{F}'$}\;
\caption{\textsc{SpacePartition}: maximal-cuboid decomposition}
\label{alg:partition}
\end{algorithm}

\noindent
\textsc{MaximalSplit}$(s, c)$ generates at most six non-overlapping
maximal sub-cuboids of $s$ after removing the intersection with $c$,
corresponding to the six possible orthogonal faces of $s$ that extend
beyond $c$:
\begin{align}
\label{eq:maxsplit}
&\text{Left:}   && s^x < c^x \;\Rightarrow\; \mathrm{Cuboid}(s^x,\, s^y,\, s^z,\,
                    c^x - s^x,\, s^H,\, s^D),                         \nonumber\\
&\text{Right:}  && s^x + s^W > c^x + c^W \;\Rightarrow\; \mathrm{Cuboid}(c^x+c^W,\, s^y,\, s^z,\,
                    s^x+s^W - (c^x+c^W),\, s^H,\, s^D),               \nonumber\\
&\text{Bottom:} && s^y < c^y \;\Rightarrow\; \mathrm{Cuboid}(s^x,\, s^y,\, s^z,\,
                    s^W,\, c^y - s^y,\, s^D),                          \nonumber\\
&\text{Top:}    && s^y + s^H > c^y + c^H \;\Rightarrow\; \mathrm{Cuboid}(s^x,\, c^y+c^H,\, s^z,\,
                    s^W,\, s^y+s^H-(c^y+c^H),\, s^D),                  \nonumber\\
&\text{Back:}   && s^z < c^z \;\Rightarrow\; \mathrm{Cuboid}(s^x,\, s^y,\, s^z,\,
                    s^W,\, s^H,\, c^z - s^z),                          \nonumber\\
&\text{Front:}  && s^z + s^D > c^z + c^D \;\Rightarrow\; \mathrm{Cuboid}(s^x,\, s^y,\, c^z+c^D,\,
                    s^W,\, s^H,\, s^z+s^D-(c^z+c^D)).                  \nonumber
\end{align}


% ============================================================
%  SECTION 8.X — STATISTICAL ANALYSIS OF RESULTS
%  Add inside the Results & Discussion section
% ============================================================

\subsection{Statistical Significance of Heuristic Improvements}
\label{sec:statistics}

To assess whether the observed improvements in volumetric utilisation are
statistically significant, a non-parametric Wilcoxon signed-rank test
\citep{Wilcoxon1945} is applied to the paired differences between each
heuristic's per-flight utilisation and the historical baseline across
the 16 evaluated flights.
The Wilcoxon test is preferred over a paired $t$-test because it makes
no normality assumption on the differences, which is important given the
small sample size ($n=16$) and the bounded nature of the utilisation
metric.

The null hypothesis $H_0$ states that the heuristic-based DSM produces
the same median volumetric utilisation as the historical manual baseline;
the alternative hypothesis $H_1$ is one-sided (heuristic $>$ baseline).
Results are reported together with descriptive statistics
(mean, standard deviation, minimum, and maximum utilisation per method)
in Table~\ref{tab:stats}.
A significance level of $p < 0.05$ is used throughout.

% NOTE TO AUTHORS: populate Table 6 (Comparative performance) with
% mean ± std values computed from the 16-flight dataset and report
% p-values from scipy.stats.wilcoxon(alternative='greater').
% See notebook 01_heuristics.ipynb for the raw per-flight results.


% ============================================================
%  SECTION 9 — CONTRIBUTIONS STATEMENT
%  Add at the END of the Introduction, before the road-map paragraph
% ============================================================

% Replace the existing road-map paragraph with the following block:

\noindent
The main contributions of this study are:
\begin{enumerate}
  \item \textbf{Aircraft-scale online 3D bin packing model.}
        We extend the DeepPack3D framework~\citep{Tsang2025} from
        single-bin evaluation to a realistic, heterogeneous multi-ULD
        configuration matching the actual cargo layout of a Boeing
        747-400F, enabling aircraft-level utilisation measurement for
        the first time in this domain.
  \item \textbf{Industry-specific problem formulation.}
        We formalise the Air Cargo Loading Problem (ACLP) for the
        cut-flower industry, explicitly modelling the online, irreversible
        nature of conveyor-based palletisation, the 50\,\% base-support
        stability constraint, and the dual-region curvature approximation
        for fuselage-adjacent containers.
  \item \textbf{Comparative evaluation on real operational data.}
        The DSM is validated on historical data from 16 full Boeing
        747-400F charter flights, demonstrating that deterministic
        heuristics — in particular the Bottom-Left policy — consistently
        outperform historical manual operations by 3--12 percentage
        points in volumetric utilisation, with statistical significance
        confirmed via a Wilcoxon signed-rank test.
  \item \textbf{Quantified operational and environmental impact.}
        We translate volumetric gains into concrete operational metrics,
        showing that the proposed DSM can recover 470--1,514 additional
        boxes per flight, reduce unit freight costs by approximately 8\%,
        and lower the CO\textsubscript{2} intensity of each shipment.
\end{enumerate}


% ============================================================
%  SECTION 9.X — LIMITATIONS
%  Add as a standalone subsection before Future Work
% ============================================================

\subsection{Limitations}
\label{sec:limitations}

Several limitations of the present study should be acknowledged.
First, the model is validated on a dataset of 16 charter flights from
a single sponsor company.
While these flights represent real worst-case operational conditions,
generalisation to other carriers, routes, or aircraft types has not been
tested.
Second, all ULDs are modelled as idealized rectangular parallelepipeds.
Although the curvature of fuselage-adjacent containers is partially
captured via a dual-region decomposition, the current implementation
does not compute precise curved-surface constraints, which may lead to
a slight overestimation of usable volume in lower-deck positions.
Third, the model treats box placement as purely volume-driven;
operational constraints such as flower fragility, required airflow
clearances, and customer-specific loading priorities are not explicitly
encoded.
Finally, only heuristic-based placement policies are evaluated in this
study; direct comparison with reinforcement learning--based agents
available in the DeepPack3D framework is reserved for future work.


% ============================================================
%  MDPI MANDATORY SECTIONS
%  Add after Conclusions, before References
% ============================================================

\subsection*{Author Contributions}
Conceptualisation, J.P., J.Z.-P.\ and R.K.;
methodology, J.P., J.Z.-P.\ and R.K.;
software, J.P.\ and J.Z.-P.;
formal analysis, J.P.\ and R.K.;
data curation, J.P.;
writing---original draft preparation, J.P., J.Z.-P.\ and R.K.;
writing---review and editing, E.G.-F., H.M.G.-A.\ and J.C.P.-P.;
supervision, E.G.-F., H.M.G.-A.\ and J.C.P.-P.;
project administration, E.G.-F.
All authors have read and agreed to the published version of the
manuscript.

\subsection*{Funding}
This research received no external funding.
% NOTE: update if any grant/fellowship applies.

\subsection*{Institutional Review Board Statement}
Not applicable.

\subsection*{Informed Consent Statement}
Not applicable.

\subsection*{Data Availability Statement}
The historical flight dataset used in this study was provided by the
sponsor company under a confidentiality agreement and is not publicly
available.
The implementation code for the adapted DeepPack3D framework is
available at \url{[REPOSITORY URL TO BE ADDED]}.

\subsection*{Conflicts of Interest}
The authors declare no conflicts of interest.
The funders had no role in the design of the study; in the collection,
analyses, or interpretation of data; in the writing of the manuscript;
or in the decision to publish the results.


% ============================================================
%  ADDITIONAL REFERENCES (add to bibliography)
% ============================================================
% The following entries should be added to the reference list.
% They appear cited in the new sections above.

% @article{Wilcoxon1945,
%   author  = {Wilcoxon, Frank},
%   title   = {Individual comparisons by ranking methods},
%   journal = {Biometrics Bulletin},
%   year    = {1945},
%   volume  = {1},
%   number  = {6},
%   pages   = {80--83},
%   doi     = {10.2307/3001968}
% }

% NOTE: Chan et al. (2006a) and (2006b) are duplicate entries —
% remove (2006b) and update all in-text citations to use (2006a) only.

% NOTE: replace informal sources with academic/institutional equivalents:
%   - "Conquerornetwork (2025)" → cite IATA or USDA floriculture report
%   - "Floriculture Market (2025)" → cite peer-reviewed market analysis
%   - "Sinchire (2025)" (newspaper) → acceptable as supplementary source;
%     add "(accessed DD Month 2025)" to the URL
